% Generated draft: requires mathematical and source review.
\documentclass[11pt]{article}
\usepackage[margin=25mm]{geometry}
\usepackage{fontspec}
\setmainfont{Latin Modern Roman}
\usepackage{amsmath,amssymb,mathtools,unicode-math}
\setmathfont{Latin Modern Math}
\usepackage{enumitem,needspace,xurl,hyperref}
\setlength{\parindent}{0pt}
\setlength{\parskip}{5pt}
\setlength{\emergencystretch}{3em}
\newcommand{\N}{\mathbb{N}}
\newcommand{\Z}{\mathbb{Z}}
\newcommand{\Q}{\mathbb{Q}}
\newcommand{\R}{\mathbb{R}}
\newcommand{\C}{\mathbb{C}}
\newcommand{\F}{\mathbb{F}}
\newcommand{\A}{\mathbb{A}}
\newcommand{\PP}{\mathbb{P}}
\newcommand{\E}{\mathbb{E}}
\newcommand{\eps}{\varepsilon}
\newcommand{\dd}{\,\mathrm{d}}
\begin{document}
\section*{UnsolvedMath ID: 30001843}
\textbf{Source problem code:} catalogue entry
\subsection{Definitions and mathematical statement}
Let \(n\) denote the integer weight of a higher polylogarithm, with particular emphasis on \(n>3\), and write \(B_n(F)\) for the higher Bloch group intended by the source. The source does not specify the admissible base fields \(F\), the definition of \(B_n(F)\), or the normalization of the polylogarithms. A functional equation here is an identity among polylogarithm values; the precise admissible arguments, coefficients, and interpretation of these identities remain unspecified.

\textbf{Problem (source-level formulation).} Construct functional equations for higher polylogarithms, especially in weights \(n>3\), that generate the universal relations needed to define \(B_n(F)\), and prove the corresponding form of Zagier's conjecture.

The phrase ``generate the universal relations'' requires a specified relation system and specified operations for deriving relations. Neither is supplied. Likewise, the corresponding form of Zagier's conjecture is not stated. Consequently, this is a research question with unresolved formulation data, rather than a fully specified quantified conjecture. In particular, the source does not determine whether one construction is required uniformly for every integer \(n>3\), or whether separate constructions in individual weights are intended.
\subsection{Short English statement}
Find functional equations for higher polylogarithms, especially for integer weights \(n>3\), that generate the relations defining the intended higher Bloch groups \(B_n(F)\), and establish the intended version of Zagier's conjecture. The field domain, group conventions, meaning of generation, and precise conjectural conclusion require clarification.
\subsection{Sources}
\begin{itemize}
\item Supplied local UnsolvedMath registry record: permanent ID \(30001843\), source problem code \texttt{catalogue entry}, title ``Functional Equations for Higher Bloch Groups,'' in the Oberwolfach Reports Open Problems collection. The record gives proposed year \(2011\) and names no proposer.
\item Supplied source citation: \emph{Explicit Methods in Number Theory (2012)}, Oberwolfach Reports, pp.\ 1971--2044, DOI \texttt{10.4171/owr/2011/35}. Supplied source and literature-source link, labelled ``original OWR report'': \url{https://doi.org/10.4171/owr/2011/35}. The background describes these reports as documenting mathematical workshops and their open-problem discussions.
\item Supplied literature-source link, labelled ``functional equations of higher logarithms'': \url{https://arxiv.org/abs/math/0207222}. No authors or additional bibliographic details for this reference are supplied.
\item Supplied literature assessment, checked on \texttt{2026-08-21}: the record labels the problem open and states that Goncharov's higher-weight conjectural programme supplies the \(n=4\) target, while no general candidate generating relation for every \(n>3\) has been established. This is supplied dated curation derived from the OWR material, not an independently verified literature assessment. The background explicitly requests independent verification of time-sensitive status and cited claims.
\item The supplied statement-status label is \texttt{corrected\_verified}. Its accompanying verification note claims that the mathematical content was checked against the official Oberwolfach report and that notation, context, and question form were normalized without changing the source-level claim. This is a supplied verification claim, not verification performed for this draft.
\end{itemize}
\end{document}
